Hager ODD (2015):

Overview:
The purpose of our model focuses on the alteration of the urban mobility system when newly-built residential or
industrial areas or rehabilitated areas are connected to the transport system. We expect it to be a beneficial tool for
researchers and urban designers to forecast mobility demands for urban areas.
Two types of entities exist in our model: the travelers in the city of Stuttgart (e.g. inhabitants or commuters) and
vehicles. Socio-demographic variables characterize the travelers, whereas frequency and capacity define the vehi-
cles. Exogenous factors for the entities are urban structure, ground plans, location of public transport stops as well as
their sizing. The City of Stuttgart is the model’s spatial extent, the temporal resolution accounts for one day with
time steps of one second. Travelers will utilize the existing circulation areas.
Every traveler pursues his personal activity chain (Fig. 1) and chooses between the existing traffic modes to trav-
el on the existing circulation areas. Public transport uses time schedules on their respective areas.
Design Concepts:
The theoretical background consists of urban morphology and population growth and their respective effects on
the urban mobility system. Agents choose their daily activity chain based on a utility function which is implemented
in MATSim (Balmer et al. 2009)
Individual decision-making consists of route, time, location and mode choice. The decision-making basis is the
utility function, afterwards the agents move based on a traffic flow simulation (Flötteröd et al. 2011). Agents com-
pete with each other for space-time-slots on the circulation areas to pursue their daily activity chain, either utilizing
empirical data or discrete-choice modelling (Horni et al. 2015). Agents observe each other and the environment.
Activity plans with the highest score in the utility function are chosen (Nagel et al. 2015).
Agents react to traffic situations and choose to drop certain activities in their daily plan after recognizing they
will be too late at the location of their next activity (replanning). Agents utilize their newly-acquired knowledge for
future activity plans.
The individual sensing of agents consists of the grasp and utilization of each traffic mode. Agents know the loca-
tion of stopping points and their travel choices affect the movement of other agents. When agents choose to alter
their activity chain, it might guide to worse results (time and/or length) than before. The traffic model is based on
the first-in-first-out-principle (Flötteröd et al. 2011).
As soon as agents predict are not able to stick with their chosen plans, they search for alternatives. In that case,
plans with the next highest score will be used (Nagel et al. 2015). The results of this process may be erroneous.
Agents perceive their environment in different ways. Interaction between agents and additional information (traf-
fic news) is possible (Dobler 2009).
Trains, busses and cars with several passengers are collectives. Agents are heterogeneous based on socio-
demographic data and their daily activity chain. The utility function consists of several variables, who differ be-
tween agents, thus leading to heterogeneous decision-making.
Initialization of the agents and the assignment of the daily activity chain is partly random.
The observation of the model’s results conduct to several different possible analyses: variances in the modal split
of the research area, congestion effects, distribution of traffic flows. Scenario-analysis are the calculation method for
this model.
Details:
The goal is to set up an agent-based model using MATSim for the city of Stuttgart. The initialization utilizes the
current modal split values of the city of Stuttgart and endures one day. The model’s input data is listed in Table 1.
The modular structure of MATSim permits the utilization of additional modules. The following modules will be
included in our agent-based model: within-day replanning, public transport, multi-modal contribution, parking,
signals and lanes, cadyts and matrix-based pt router (Horni and Nagel 2015).
In summary, the functionalities of the existing MATSim version cover most of the required features for this
study. Fig. 1 represents the proposed research design and its anticipated fields of impact. Following this, we present-
ed a concept for agent-based models. Before finally setting up the model, demographic data need to be utilized to
validate the hypothesis of this study that the City of Stuttgart experienced a significant amount of population growth
during the last years. First results regarding the demographic status of the City of Stuttgart are presented.


The basis of an agent-based MATSim simulation is a synthetic population of the case study area. For Switzerland, a range of rich
data sources are available that make it possible to set up such a population with detailed socio-demographic and behavioral attributes.
In the following, the process will be described in brief. A more detailed description of the Swiss data sets and how they have been
processed can be found in (H¨orl, 2020). Furthermore, (H¨orl and Balac, 2020) describes the methodology in detail for the case of Paris
and ˆIle-de-France. There, unlike Switzerland, all necessary data sets are publicly and openly available to replicate the process.
Initially, the synthetic Swiss population is based on the STATPOP data set, which is a comprehensive census of Switzerland,
including attributes like age and gender of all residents, but - most notably - also location by coordinate of each household. Persons and
households from this data set are used directly as a basis of the Swiss synthetic population. Afterward, work or educational locations for
each synthetic person are added. For that, the Structural Survey is used, which, unlike the registration-based census data, is a detailed
household survey conducted for around 3% of the Swiss population each year. From this data set, an origin–destination matrix between
home and work/education municipalities in Switzerland can be derived, then used to proportionately sample destination zones for
each agent, based on their home zone. In the third step, the national Mobility and Transport Microcensus is employed: a classical
household travel survey providing full day activity patterns for around 60,000 respondents. We use a statistical matching procedure to
sample and attach activity chains to the synthetic persons, based on socio-demographic attributes. This way, activity patterns
correlating with certain age groups or gender are assigned according to the population. These activity chains contain activity types,
activity durations and connecting modes of transport as reported in the survey data. As departure and arrival times are rounded
strongly (in five minute intervals) as is the case in many surveys, activity chains (including all departure and arrival times) are shifted
by a random offset of +/�� 30 min for each person to avoid having unrealistically many people departing at the same time at one
location, which would, in turn, affect the traffic dynamics in the simulation negatively. Finally, specific work and education locations
are sampled, based on the destination zones of the agents, from STATENT, the national enterprise census, which lists all facilities in
Switzerland by coordinate. Further, secondary activities in the agents’ plans (such as shopping or leisure), are assigned to discrete
facilities using a specifically-designed algorithm for secondary location assignment (H¨orl and Axhausen, 2020). To complete the
simulation data, an exhaustive digital public transit schedule for Switzerland is integrated into the simulation and the road network is
extracted from OpenStreetMap, including all roads from motorways to residential roads. In total, this process yields a synthetic
representation of Switzerland that reproduces the daily mobility needs of the population and available infrastructure.


